% !TeX root = ../main.tex

% ==================================================
% A.2. 중첩 함수와 Lexical Scope
% ==================================================

\subsection{중첩 함수와 Lexical Scope}

Nested Functions · Lexical Scoping · Function Definition, Execution

\subsubsection{이전 내용과의 연결}

\zsh{A.1}에서는 Python의 함수가 단순한 코드 블록이 아니라

\zsh{실제로 생성되고 참조되는 함수 객체(Function Object)}라는 사실을 학습했다.

특히 다음 네 가지를 구분하였다.

\begin{dhksmry}
    1. def 문장의 실행은 함수 객체를 생성하고 이름을 바인딩한다.

    2. func은 함수 객체를 참조하지만, func()은 그 객체를 호출한다.

    3. 함수 객체는 다른 함수의 인수로 전달하거나 반환할 수 있다.

    4. 함수 호출 시 인수 객체가 매개변수에 바인딩되며, 
    
       return은 호출 결과를 반환한다.
\end{dhksmry}

이번에는 이 원리들을 함수 내부에서 또 다른 함수를 정의하는 경우로 확장할 것이다.

다음 코드를 살펴보자.

\begin{dhkcode}
def outer():
    x = 10

    def inner():
        return x

    return inner

f = outer()

print(f())
\end{dhkcode}

예상 실행 결과:

\begin{dhkzsh}
10
\end{dhkzsh}

여기에는 중요한 질문이 숨어 있다.

\zsh{outer()}의 호출은 이미 끝났는데, 왜 나중에 호출한 \zsh{inner()}는 \zsh{outer()} 내부에서 정의한 \zsh{x}에 접근할 수 있을까?

이 현상을 설명하려면 다음 개념을 구분해야 한다.

\begin{dhksmry}
    - 함수가 어디에서 정의되었는가?

    - 함수 객체는 언제 생성되는가?

    - 함수 본문은 언제 실행되는가?

    - 함수 내부의 이름은 어떤 스코프를 기준으로 탐색되는가?

    - 외부 함수의 실행이 종료되면 그 지역 변수의 바인딩은 어떻게 되는가?
\end{dhksmry}

마지막 질문은 \zsh{A.4}의 클로저와 직접 연결된다.

이번 \zsh{A.2}에서는 그에 앞서 중첩 함수의 정의 구조와 Lexical Scope의 원리를 확립할 것이다.

\subsubsection{중첩 함수(Nested Function)의 정의}

중첩 함수란 다른 함수의 본문 내부에서 정의되는 함수를 말한다.

가장 기본적인 구조는 다음과 같다.

\begin{dhkcode}
def outer():

    def inner():
        return "Hello"

    return inner()
\end{dhkcode}

여기서 \zsh{outer}는 외부 함수(Outer Function), \zsh{inner}는 내부 함수(Inner Function)이다.

하지만 이 용어가 두 함수 사이에 특별한 상속 관계가 있다는 의미는 아니다.

\zsh{outer}와 \zsh{inner}는 각각 독립적인 함수 객체로 생성될 수 있다.

두 함수 사이의 중요한 관계는 정의된 위치에 따른 렉시컬 스코프의 중첩이다.

\begin{dhknote}
\begin{dhksmry}[centered]
중첩 함수의 기본 구조
\end{dhksmry}
\end{dhknote}

\begin{dhkdir}[in=6]
모듈 스코프(Module Scope)
> outer 함수 정의
>> outer 함수의 본문
>>> 지역 변수 x
>>> inner 함수 정의
>>>> inner 함수의 본문
\end{dhkdir}

여기서 주의할 점이 있다.

위 그림은 함수의 소스 코드상 정의 위치를 나타내는 것이지, 함수 객체가 모두 동시에 생성되어 있다는 뜻은 아니다.

이 둘을 구별하는 것이 이번 단원의 첫 번째 핵심이다.

\subsubsection{중첩 함수의 생성 시점}

다음 코드를 살펴보자.

\begin{dhkcode}
def outer():
    print("outer 실행 시작")

    def inner():
        print("inner 실행")
        return 10

    print("inner 정의 완료")

    return inner

print("모듈 코드 실행")

f = outer()

print("outer 호출 완료")

result = f()

print(result)
\end{dhkcode}

예상 실행 결과:

\begin{dhkzsh}
모듈 코드 실행
outer 실행 시작
inner 정의 완료
outer 호출 완료
inner 실행
10
\end{dhkzsh}

이 결과를 정확히 이해하려면 함수 객체 생성과 함수 본문 실행을 순서대로 추적해야 한다.

\begin{dhknote}
\begin{dhksmry}[centered]
1. outer 함수 객체 생성
\end{dhksmry}
\end{dhknote}

Python이 다음 정의문을 실행한다.

\begin{dhkcode}
def outer():
    ...
\end{dhkcode}

이때 \zsh{outer} 함수 객체가 생성되고 모듈 네임스페이스의 \zsh{outer}라는 이름에 바인딩된다.

그러나 \zsh{outer} 본문은 실행되지 않는다.

따라서 이 시점에는 \zsh{outer} 본문에 있는 \zsh{def inner} 문장도 실행되지 않았다.

여기서 중요한 구분이 있다.

Python은 소스 코드를 컴파일하는 과정에서 중첩 함수의 코드 객체를 준비할 수 있다.

하지만 그것이 곧 실행 시점의 \zsh{inner} 함수 객체가 생성되었다는 의미는 아니다.

코드 객체(Code Object)와 함수 객체(Function Object)는 서로 다른 객체다.

\begin{dhknote}
\begin{dhksmry}[centered]
2. outer() 호출
\end{dhksmry}
\end{dhknote}

\begin{dhkcode}
f = outer()
\end{dhkcode}

이제 \zsh{outer} 함수 객체가 호출된다.

일반적인 Python 함수 호출에서는 해당 호출을 위한 실행 프레임(Frame)이 생성되고, 매개변수와 지역 이름을 관리하는 실행 환경이 마련된다.

이 프레임에서 다음 문장이 실행된다.

\begin{dhkcode}
print("outer 실행 시작")
\end{dhkcode}

첫 번째 함수 내부 출력이 발생한다.

\begin{dhknote}
\begin{dhksmry}[centered]
3. inner 함수 객체 생성
\end{dhksmry}
\end{dhknote}

이어서 다음 정의문에 도달한다.

\begin{dhkcode}
def inner():
    print("inner 실행")
    return 10
\end{dhkcode}

이때 새로운 함수 객체가 생성된다.

그리고 \zsh{outer()} 호출의 지역 스코프에서 \zsh{inner}라는 이름이 그 객체에 바인딩된다.

그러나 \zsh{inner} 함수 본문은 아직 실행되지 않았다.

즉, 현재 상태는 다음과 같다.

\begin{dhkdir}[in=6]
모듈 네임스페이스
> outer -> outer 함수 객체
\end{dhkdir}

\begin{dhkdir}[in=6]
outer() 실행 환경
> inner -> inner 함수 객체
\end{dhkdir}

여기서 두 함수 객체는 서로 다른 객체다.

\begin{dhkcode}
type(outer)
type(inner)
\end{dhkcode}

두 객체 모두 일반적인 사용자 정의 함수라면 타입은 \zsh{function}이지만, \zsh{inner}라는 지역 이름은 현재 \zsh{outer()} 실행 환경에서만 직접 접근할 수 있다.

따라서 모듈 수준에서 곧바로 \zsh{inner}를 참조하려고 하면 

일반적으로 \zsh{NameError}가 발생한다.

\begin{dhknote}
\begin{dhksmry}[centered]
4. inner 함수 객체 반환
\end{dhksmry}
\end{dhknote}

\begin{dhkcode}
return inner
\end{dhkcode}

여기서 \zsh{inner()}가 아니라 \zsh{inner}를 반환한다는 점에 주목해야 한다.

\zsh{inner}가 참조하는 함수 객체 자체를 반환하는 것이다.

따라서 다음 문장의 실행이 완료되면

\begin{dhkcode}
f = outer()
\end{dhkcode}

모듈 네임스페이스의 \zsh{f}가 반환된 함수 객체를 참조하게 된다.

\begin{dhkdir}[in=6]
모듈 네임스페이스
> outer -> outer 함수 객체
> f -> inner 함수 객체
\end{dhkdir}

이후 \zsh{outer()} 호출은 종료된다.

그러나 반환된 \zsh{inner} 함수 객체는 \zsh{f}가 참조하고 있으므로 계속 존재할 수 있다.

\begin{dhknote}
\begin{dhksmry}[centered]
5. 반환된 함수 객체 호출
\end{dhksmry}
\end{dhknote}

마지막으로 다음 문장이 실행된다.

\begin{dhkcode}
result = f()
\end{dhkcode}

\zsh{f}는 \zsh{inner} 함수 객체를 참조하고 있다.

따라서 \zsh{inner} 함수 본문이 실행되고 10을 반환한다.

이 결과가 \zsh{result}에 바인딩되는 것이다.

\begin{dhknote}
\begin{dhksmry}[centered]
핵심 구분
\end{dhksmry}
\end{dhknote}

\begin{dhktable}
|사건|	발생 시점|
|:---|:---|
|outer 함수 객체 생성|	모듈에서 def outer 실행 시|
|inner 함수 객체 생성|	outer() 실행 중 def inner에 도달할 때|
|outer 함수 본문 실행|	outer() 호출 시|
|inner 함수 본문 실행|	inner() 또는 f() 호출 시|
|inner 함수 객체 반환|	outer()의 return inner 실행 시|
\end{dhktable}

여기까지가 중첩 함수의 기본 실행 구조다.

\subsubsection{중첩 함수는 호출할 때마다 새로 생성되는가?}

이제 다음 코드를 살펴보자.

\begin{dhkcode}
def outer():

    def inner():
        return 10

    return inner


f1 = outer()
f2 = outer()

print(f1 is f2)

print(f1.__code__ is f2.__code__)
\end{dhkcode}

일반적인 CPython 실행 결과:

\begin{dhkzsh}
False
True
\end{dhkzsh}

왜 이런 차이가 발생할까?

\begin{dhknote}
\begin{dhksmry}[centered]
함수 객체의 식별성
\end{dhksmry}
\end{dhknote}

먼저 다음 두 호출을 살펴보자.

\begin{dhkcode}
f1 = outer()
f2 = outer()
\end{dhkcode}

첫 번째 \zsh{outer()} 호출에서 \zsh{def inner} 문장이 실행되면서 함수 객체가 생성된다.

두 번째 \zsh{outer()} 호출에서도 같은 정의문이 다시 실행되므로 새로운 함수 객체가 생성된다.

따라서 두 함수 객체는 동일한 객체가 아니다.

\begin{dhkcode}
f1 is f2
\end{dhkcode}

결과는 \zsh{False}가 된다.

그러데 두 함수의 코드 객체는 어떨까?

\begin{dhkcode}
f1.__code__ is f2.__code__
\end{dhkcode}

일반적인 CPython에서는 두 함수 객체가 동일한 코드 객체를 참조한다.

왜냐하면 같은 함수 정의에 대응하는 컴파일된 코드 객체를 재사용하면서 서로 다른 함수 객체를 생성할 수 있기 때문이다.

다만 코드 객체의 \zsh{is} 동일성은 Python 언어 차원에서 모든 구현에 보장되는 일반 규칙으로 취급하면 안 된다.

여기서 확실하게 구분해야 하는 것은 다음이다.

동일한 함수 정의에서 생성된 함수 객체라고 해서 반드시 동일한 함수 객체인 것은 아니다.

반면 같은 함수 객체를 여러 이름에 바인딩하는 것은 새로운 함수 객체 생성과 다르다.

\begin{dhkcode}
f1 = outer()
f2 = f1

print(f1 is f2)
\end{dhkcode}

출력:

\begin{dhkzsh}
True
\end{dhkzsh}

이 경우에는 함수 객체를 새로 생성하지 않고 기존 객체에 이름을 추가로 바인딩했기 때문이다.

이 차이는 \zsh{A.4}에서 여러 클로저 인스턴스가 서로 독립적인 상태를 유지하는 원리를 이해할 때 중요해진다.

\subsubsection{Lexical Scope란 무엇인가?}

이제 이번 학습의 중심 개념으로 넘어가 보자.

\zsh{Lexical Scope(렉시컬 스코프)}란 함수 내부에서 사용한 이름이 어떤 외부 스코프와 연결되는지가 소스 코드상 함수의 정의 위치에 따라 결정되는 규칙이다.

이를 정적 스코프(Static Scope)라고도 부른다.

여기서 중요한 것은 함수가 어디에서 호출되었는가가 아니라 어디에서 정의되었는가이다.

다음 코드를 보자.

\begin{dhkcode}
x = "Global"

def outer():
    x = "Outer"

    def inner():
        print(x)

    return inner

def caller(func):
    x = "Caller"
    func()

f = outer()

caller(f)
\end{dhkcode}

예상 실행 결과:

\begin{dhkzsh}
Outer
\end{dhkzsh}

이 결과는 매우 중요하다.

왜 \zsh{"Caller"}가 아니라 \zsh{"Outer"}가 출력될까?

\begin{dhknote}
\begin{dhksmry}[centered]
실행 과정 분석
\end{dhksmry}
\end{dhknote}

먼저 모듈 스코프에 다음 이름이 바인딩된다.

\begin{dhkcode}
x = "Global"
\end{dhkcode}

그다음 \zsh{outer}와 \zsh{caller} 함수 객체가 생성된다.

이후 다음 표현식이 실행된다.

\begin{dhkcode}
f = outer()
\end{dhkcode}

\zsh{outer()} 내부에서는 지역 이름 \zsh{x}가 문자열 \zsh{"Outer"}를 참조하게 된다.

그리고 \zsh{inner} 함수 객체가 생성된ㄷ.

이때 \zsh{inner} 함수는 소스 코드상 \zsh{outer} 내부에 정의되어 있다.

따라서 \zsh{inner}의 렉시컬 외부 스코프는 \zsh{outer}의 지역 스코프다.

이후 \zsh{inner} 함수 객체가 반환되어 \zsh{f}에 바인딩된다.

이제 다음 호출을 살펴보자.

\begin{dhkcode}
caller(f)
\end{dhkcode}

\zsh{caller()}의 매개변수 \zsh{func}는 \zsh{f}가 참조하는 \zsh{inner} 함수 객체를 참조하게 된다.

\zsh{caller()} 내부에서는 다음 문장이 실행된다.

\begin{dhkcode}
x = "Caller"
\end{dhkcode}

그러나 이 \zsh{x}는 \zsh{caller()} 호출의 지역 이름이다.

\zsh{inner} 함수가 정의된 렉시컬 구조에는 \zsh{caller}의 지역 스코프가 포함되어 있지 않다.

따라서 다음 호출에서

\begin{dhkcode}
func()
\end{dhkcode}

\zsh{inner}가 실행되더라도 \zsh{caller}의 지역 변수 \zsh{x}를 탐색하지 않는다.

\zsh{inner}는 자신이 정의된 외부 함수 \zsh{outer}의 바인딩을 사용한다.

그래서 \zsh{"Outer"}가 출력되는 것이다.

\begin{dhknote}
\begin{dhksmry}[centered]
Lexical Scope와 Dynamic Scope 비교
\end{dhksmry}
\end{dhknote}

\begin{dhktable}
|구분|	Lexical Scope|	Dynamic Scope|
|:---|:---|:---|
|이름 탐색의 기준|	함수가 정의된 코드 구조|	함수 호출의 동적 경로|
|외부 이름의 연결|	렉시컬 외부 스코프|	호출자 측 실행 환경|
|Python의 일반적인 함수 이름 탐색|	사용|	사용하지 않음|
|위 예제의 x|	"Outer"|	가상의 동적 스코프라면 "Caller"가 될 수 있음|
\end{dhktable}

Python은 일반적인 함수 이름 탐색에 Lexical Scope를 사용한다.

따라서 함수 객체를 다른 함수에 전달하거나 다른 위치에서 호출한다고 해서 그 함수의 렉시컬 외부 스코프가 변경되지는 않는다.

이것이 Lexical Scope의 핵심이다.

\subsubsection{정의 위치와 호출 위치는 왜 구분해야 할까?}

앞의 예제를 조금 변형해 보자.

\begin{dhkcode}
x = 100

def outer():
    x = 200

    def inner():
        return x

    return inner

f = outer()

def execute():
    x = 300
    return f()

print(execute())
\end{dhkcode}

예상 실행 결과:

\begin{dhkzsh}
200
\end{dhkzsh}

여기서 \zsh{inner} 함수는 \zsh{execute()} 내부에서 호출되고 있다.

하지만 \zsh{execute()}의 지역 변수 \zsh{x=300}은 \zsh{inner}의 이름 탐색에 영향을 주지 않는다.

또한 모듈 스코프에는 \zsh{x=100}이 존재하지만, \zsh{inner}는 자신을 감싸는 \zsh{outer}의 \zsh{x}를 사용한다.

이 관계를 렉시컬 구조로 나타내면 다음과 같다.

\begin{dhkdir}[in=6]
Global Scope
> x = 100
> outer
>> x = 200
>> inner
>>> return x
> execute
>> x = 300
>> f()
\end{dhkdir}

\zsh{inner}가 참조하는 \zsh{x}의 탐색 관계는 다음과 같다.

\begin{dhksmry}[centered]
inner의 Local
\da
outer의 Enclosing
\da
Global
\da
Builtins
\end{dhksmry}

여기서는 \zsh{outer}의 \zsh{x}를 찾으므로 그 값을 사용한다.

반면 \zsh{execute}는 \zsh{inner}의 렉시컬 외부 스코프에 포함되지 않는다.

함수 호출 관계(Call Relationship)와 렉시컬 스코프 관계(Lexical Scope Relationship)는 서로 다른 구조라는 사실을 기억해야 한다.

이것은 이후 클로저와 데코레이터에서 매우 중요해진다.

\subsubsection{중첩 함수는 외부 함수의 모든 지역 변수를 사용할 수 있는가?}

중첩 함수는 자신을 감싸는 외부 함수의 지역 이름을 참조할 수 있다.

하지만 모든 이름이 자동으로 클로저에 저장되는 것은 아니다.

다음 코드를 살펴보자.

\begin{dhkcode}
def outer():
    x = 10
    y = 20

    def inner():
        return x

    return inner

f = outer()

print(f())
print(f.__code__.co_freevars)
\end{dhkcode}

예상 실행 결과:

\begin{dhkzsh}
10
('x',)
\end{dhkzsh}

여기서 \zsh{inner}는 \zsh{x}만 참조하고 있다.

\zsh{y}는 \zsh{outer}의 지역 변수이지만 \zsh{inner}에서 사용되지 않는다.

따라서 \zsh{inner}의 자유 변수(Free Variable) 목록에는 \zsh{x}만 나타난다.

여기서 자유 변수란 현재 함수의 지역 스코프에서 바인딩되지 않았지만, 렉시컬 외부 함수 스코프의 바인딩을 사용하는 이름을 말한다.

\zsh{co_freevars}는 코드 객체가 어떤 자유 변수 이름을 사용하는지 확인하는 데 유용하다.

다만 지금은 자유 변수가 존재한다는 사실만 확인하면 충분하다.

자유 변수가 실제로 어떤 셀 객체와 연결되고, 외부 함수 호출이 종료된 뒤에도 어떻게 유지되는지는 \zsh{A.4}에서 자세히 분석할 것이다.

\subsubsection{외부 함수의 지역 변수는 언제 결정되는가?}

Lexical Scope를 학습할 때 자주 발생하는 오해가 있다.

함수의 정의 위치가 이름 탐색의 스코프를 결정한다는 사실과, 그 이름에 바인딩된 객체가 함수 정의 시점에 고정된다는 것은 서로 다른 이야기다.

다음 코드를 보자.

\begin{dhkcode}
def outer():
    x = 10

    def inner():
        return x

    x = 20

    return inner

f = outer()

print(f())
\end{dhkcode}

예상 실행 결과:

\begin{dhkzsh}
20
\end{dhkzsh}

왜 10이 아니라 20일까?

\zsh{inner} 함수 객체가 생성되는 시점에는 \zsh{x}가 10을 참조하고 있다.

하지만 \zsh{inner}는 단순히 그 시점의 정수 객체 10을 복사해서 보관하는 것이 아니다.

\zsh{inner}가 사용하는 외부 이름 \zsh{x}는 외부 함수의 바인딩과 연결된다.

따라서 이후 다음 문장이 실행되면

\begin{dhkcode}
x = 20
\end{dhkcode}

외부 함수의 \zsh{x} 바인딩이 새로운 정수 객체 20을 참조하도록 변경된다.

그리고 \zsh{inner}를 호출하면 현재 연결된 바인딩을 통해 20을 얻는 것이다.

이것이 이후 클로저를 학습할 때 중요한 출발점이 된다.

\begin{dhknote}
\begin{dhksmry}[centered]
더 극단적인 사례
\end{dhksmry}
\end{dhknote}

\begin{dhkcode}
def outer():

    def inner():
        return x

    x = 100

    return inner

f = outer()

print(f())
\end{dhkcode}

예상 실행 결과:

\begin{dhkzsh}
100
\end{dhkzsh}

\zsh{inner}가 정의되는 순간에는 아직 \zsh{x = 100}이 실행되지 않았다.

그런데도 \zsh{inner()}는 정상적으로 100을 반환한다.

왜냐하면 이름의 렉시컬 연결과 실제 값의 조회 시점이 다르기 때문이다.

단, 이 예제에서 \zsh{x = 100}이 실행되기 전에 \zsh{inner()}를 호출하면 바인딩되지 않은 자유 변수를 읽으려 하므로 \zsh{NameError}가 발생할 수 있다.

이 오류는 \zsh{A.3}에서 지역 변수의 정적 판정과 함께 분석할 것이다.

\subsubsection{중첩 함수와 일반 함수의 차이}

지금까지의 내용을 비교해 보자.

\begin{dhkcode}
x = 10

def normal():
    return x

def outer():
    x = 20

    def inner():
        return x

    return inner

f = outer()

print(normal())
print(f())
\end{dhkcode}

예상 실행 결과:

\begin{dhkzsh}
10
20
\end{dhkzsh}

두 함수 모두 본문에서 \zsh{x}라는 이름을 사용한다.

하지만 \zsh{normal}은 모듈 스코프에서 정의되었고, \zsh{inner}는 \zsh{outer} 내부에서 정의되었다.

따라서 이름 탐색 경로가 다르다.

\begin{dhktable}
|구분|normal|inner|
|:---|:---|:---|
|정의 위치|모듈 스코프|outer 내부|
|x의 종류|전역 이름|자유 변수|
|외부 탐색 대상|모듈 전역 네임스페이스|outer의 렉시컬 스코프|
|반환값|10|20|
\end{dhktable}

이 차이는 단순히 중첩 함수가 특별한 문법을 사용하기 때문에 발생하는 것이 아니다.

함수가 정의된 렉시컬 환경에 따라 이름 탐색의 대상이 달라지기 때문이다.

\subsubsection{중첩 함수와 클래스 메서드의 차이}

이제 기존 클래스 객체 모델과 연결해 보자.

다음 코드를 살펴보자.

\begin{dhkcode}
class Example:
    x = 100

    def method(self):
        return x

obj = Example()
\end{dhkcode}

여기서 \zsh{method()} 내부의 \zsh{x}는 클래스 네임스페이스에 있는 \zsh{Example.x}를 자동으로 참조할까?

그렇지 않다.

일반적인 메서드의 함수 본문에서 이름을 탐색할 때 클래스 네임스페이스가 자동으로 Enclosing 함수 스코프처럼 취급되는 것은 아니다.

따라서 모듈 수준에 \zsh{x}가 정의되어 있지 않다면 다음 호출은 \zsh{NameError}를 발생시킨다.

\begin{dhkcode}
obj.method()
\end{dhkcode}

반면 다음 코드는 정상적으로 동작한다.

\begin{dhkcode}
class Example:
    x = 100

    def method(self):
        return self.x

obj = Example()

print(obj.method())
\end{dhkcode}

예상 실행 결과:

\begin{dhkzsh}
100
\end{dhkzsh}

여기서는 이름 탐색과 속성 탐색을 구분해야 한다.

\begin{dhkcode}
return x
\end{dhkcode}

는 이름 탐색(Name Resolution)이다.

반면

\begin{dhkcode}
return self.x
\end{dhkcode}

는 먼저 \zsh{self}라는 이름을 탐색한 뒤, 그 객체의 \zsh{x}속성을 탐색하는 표현식이다.

따라서 \zsh{self.x}는 기존에 배운 인스턴스·클래스 숙성 탐색 규칙에 따라 처리된다.

이것은 Lexical Scope와 객체의 속성 탐색이 서로 다른 메커니즘이라는 사실을 보여준다.

단, 클래스 본문 안에서 정의된 모든 중첩 함수가 클래스 네임스페이스를 절대 참조할 수 없다는 뜻은 아니다.

특히 클래스 본문에서 직접 실행되는 표현식, 컴프리헨션, \zsh{__class__} 셀과 관련된 특수한 규칙은 별도로 존재한다.

이번 학습에서는 일반적인 인스턴스 메서드의 이름 탐색과 \zsh{self}를 통한 속성 탐색이 다르다는 점을 확립하면 된다.

\subsubsection{실제 응용 - 함수 생성기(Function Factory)}

중첩 함수의 실용적인 활용을 살펴보자.

\begin{dhkcode}
def make_multiplier(factor):

    def multiply(value):
        return factor * value

    return multiply


double = make_multiplier(2)
triple = make_multiplier(3)

print(double(10))
print(triple(10))
\end{dhkcode}

예상 실행 결과:

\begin{dhkzsh}
20
30
\end{dhkzsh}

이 코드는 함수 객체를 생성하여 반환하는 함수 생성기(Function Factory)의 예제다.

\begin{dhknote}
\begin{dhksmry}[centered]
첫 번째 호출
\end{dhksmry}
\end{dhknote}

\begin{dhkcode}
double = make_multiplier(2)
\end{dhkcode}

\zsh{make_multiplier()}가 호출되면서 \zsh{factor}에 정수 객체 2가 바인딩된다.

그 호출 중 \zsh{multiply} 함수 객체가 생성된다.

이 함수 객체는 외부 함수의 \zsh{factor} 바인딩을 참조하는 구조를 가진다.

이후 \zsh{multiply} 함수 객체가 반환되고 \zsh{double}에 바인딩된다.

\begin{dhknote}
\begin{dhksmry}[centered]
두 번째 호출
\end{dhksmry}
\end{dhknote}

\begin{dhkcode}
triple = make_mulitplier(3)
\end{dhkcode}

이번에는 별도의 \zsh{make_muliplier()} 호출이 발생한다.

새로운 실행 환경에서 \zsh{factor}에 3이 바인딩되고, 새로운 \zsh{multiply} 함수 객체가 생성된다.

따라서 \zsh{double}과 \zsh{triple}은 서로 다른 함수 객체다.

\begin{dhkcode}
print(double is triple)
\end{dhkcode}

출력:

\begin{dhkzsh}
False
\end{dhkzsh}

각 함수는 서로 다른 외부 함수 호출에서 형성된 바인딩을 참조하므로 서로 다른 결과를 반환할 수 있다.

여기까지 이해했다면 다음 질문이 자연스럽게 발생할 것이다.

외부 함수 호출이 끝났는데도 \zsh{factor}에 대한 바인딩이 어떻게 유지되는가?

이 질문이 바로 \zsh{A.4}에서 다룰 클로저의 핵심이다.

현재는 함수가 정의된 렉시컬 환경과 실제 호출 시점이 분리될 수 있다는 사실에 집중한다.

\subsubsection{핵심 정리}

이번 학습에서 가장 중요한 것은 다음 다섯 가지다.

\begin{dhknote}
\begin{dhksmry}[centered]
1. 중첩 함수는 함수 내부에서 정의되는 함수다.
\end{dhksmry}
\end{dhknote}

\begin{dhkcode}
def outer():
    def inner():
        pass
\end{dhkcode}

두 함수는 서로 다른 함수 객체가 될 수 있다.

\begin{dhknote}
\begin{dhksmry}[centered]
2. 함수 객체 생성과 함수 본문 실행은 다르다.
\end{dhksmry}
\end{dhknote}

\zsh{outer()}가 호출되어 \zsh{def inner} 문장에 도달하면 \zsh{inner} 함수 객체가 생성된다.

그러나 \zsh{inner} 본문은 \zsh{inner()}가 호출될 때 실행된다.

\begin{dhknote}
\begin{dhksmry}[centered]
3. Lexical Scope는 정의 위치를 기준으로 한다.
\end{dhksmry}
\end{dhknote}

함수가 어디에서 호출되는지가 아니라 어디에서 정의되었는지가 렉시컬 외부 스코프를 결정한다.

\begin{dhknote}
\begin{dhksmry}[centered]
4. 이름의 탐색 대상과 바인딩된 값은 구분해야 한다.
\end{dhksmry}
\end{dhknote}

함수가 외부 이름을 참조한다고 해서 함수 정의 시점의 값을 반드시 복사하는 것은 아니다.

\begin{dhknote}
\begin{dhksmry}[centered]
5. 함수 호출 관계와 스코프 관계는 다르다.
\end{dhksmry}
\end{dhknote}

다른 함수가 중첩 함수를 호출하더라도 호출자의 지역ㄷ 스코프가 그 중첩 함수의 Enclosing Scope가 되는 것은 아니다.
